\section{Limitations}
\label{sec:limitations}

\subsection{Internal validity}
\label{sec:limits-internal}

The simulated rows are sample quantiles at 99.9 percent from $10^6$ draws, so each rests on 1,000 tail draws, and their standard
errors are about one percent of a row. The comparisons that matter are two orders of magnitude larger than this error (the width
of the band is \bdWidthOverSe{} standard errors), but the coverage of a structure whose capital lies within a few standard errors of
the edge of an interval is sensitive to it: the reflected Clayton copula is inside the interval of both matrices near its upper end,
and a different seed could move it across. The widening by three standard errors of the capital of the structure, the tolerance of
5 percent, the 75 percent of the coverage and the 100 standard errors of the width are thresholds of the author, fixed before the
experiments and not derived. The capital of a planted structure is itself an estimate, from 20 samples of $10^6$ draws, with a
standard error of at most CLP \bdTruthSeMax{} billion. The rearrangement algorithm that computes the bounds has no proof of
convergence; the evidence here is the narrowness of the ranges, and a marginal with another shape of tail could behave otherwise.
The coverage of the planted structures by the sharp range is, by Proposition~\ref{prop:coverage}, a check of the implementation and
not a finding. The classification of a test as unrunnable without a row is made by the arguments of the function for two of the
tests and by an engine call for one.

\subsection{External validity}
\label{sec:limits-external}

All results are for one synthetic bank, four risk types, one level of 99.9 percent and one horizon of one year. The marginals are
simulated by the producers of the case, and two of them rest on conventions that the case records as conventions: the
market marginal is moved from 99 percent and one day to 99.9 percent and one year by the normal quantile ratio and the square root
of 250, and the correlation of interest rate risk with the other types is set equal to that of market risk, which the case records as unverified. Each moves every row
of the band alike, so the band does not remove it; the band shows the dependence uncertainty conditional on the marginals and not
the uncertainty of the marginals. The marginals are read from a grid of $\bdGridPoints$ probabilities that ends at $5\times10^{-7}$ on each
side with linear interpolation and clamping, so the rearrangement for the worst case reads the quantile functions only as far as
the grid allows; the lower end of the grid reaches losses of CLP \bdMarketGridMin{} billion for the market marginal, and a bound row
that depends on the extreme tail beyond the grid would not see it. The ten planted structures are chosen by the author, include
several that are in the family of the band, and are not a sample of the dependence structures of banks: the proportions covered
are properties of this set. The competing single-figure estimators were calibrated on data of the planted structure, which is
favourable to them, and the band's interior rows use named conventions; with real data a bank would have a noisier matrix.

\subsection{Construct validity}
\label{sec:limits-construct}

The experiments measure coverage of the capital of planted structures with the marginals of the case, so they measure whether the
band brackets the effect of the dependence on the capital, which is the construct of the dependence band. They do not measure
whether the band is useful for a decision, whether it changes what a bank holds, or whether a supervisor would accept it. The
width is used as the price of coverage, and a more informative comparison would be a proper scoring rule for an interval, which
is not computed here. The planted structures share the marginals of the case, so a structure that changes the marginals, such as a
dependence between the frequency and the severity of operational losses, is outside what they test; the band is conditional on
the marginals by design. The Euler allocation is evaluated for the expected shortfall of the Gaussian copula row only, and the
model-risk measures are computed on the band of that one aggregate; neither is validated against a known allocation or a known
model risk.

\subsection{Scope}
\label{sec:limits-scope}

The study does not claim that the band is optimal, that no other interval construction is better, or that the band's rows are the
only ones that a validator should show; the bounds with a constraint on the variance of the sum \citep{brv2015}, for example, give
narrower intervals under extra information and are not run here. It does not claim that no maintained toolkit or published practice
reports the band beside every aggregate, since no search beyond the cited sources was made. It does not examine the aggregation of
the expected shortfall, only of the value at risk at one level, nor a number of risk types other than four. The acceptance criteria
of the tests are internal to the validation, and no result is a statement about regulatory compliance. Nor does it claim that eight
is the right number of rows: the set of degrees of freedom is a convention, and a set with another number of values gives a band
with more or fewer rows under the same rule, so the claim is about the rule and the number of rows follows from the set.
